\FloatBarrier
\section{Appendix E. External Benchmark Details}
\phantomsection\label{app:e}
\subsection{Original-task Reruns and Scoring}
We reran BERT~\citep{devlin2019bert}, SpanBERT~\citep{joshi2020spanbert}, JobBERT, and JobSpanBERT~\citep{zhang2022skillspan} separately for skills and knowledge on SkillSpan's public HOUSE and TECH subsets. The MaChAmp~\citep{vandergoot2021machamp} \texttt{seq\_bio} CRF runs used 20 epochs, development-set checkpoint selection, and five seeds listed in the \href{https://github.com/AlfredJamesLi/chinese-skillspan-benchmark/blob/36f009fab01b8433355cb7ef5af8fad02c6cc519/reproduction/process_archive_20260924/source/tex/external_benchmarks/paper_appendix.tex}{experiment records}. BIG was unavailable; multi-task training was not rerun. Table~\ref{tab:external-native} gives MaChAmp per-site span-F1. The \href{https://github.com/AlfredJamesLi/chinese-skillspan-benchmark/blob/36f009fab01b8433355cb7ef5af8fad02c6cc519/reproduction/process_archive_20260924/README.md#skillspan-pooled-scores-previous-table-25}{archived pooled scores} report BIO rescoring of the same predictions with the correct task column. These scoring procedures yield different results. The native runs contain 3,570 test sentences, compared with 3,569 in the common-protocol release; the one-record difference has not been resolved by matching IDs. For published test references, we use only the original paper's TEST/STL rows. Its HOUSE/TECH rows report development results.

For the German public ICT task~\citep{gnehm2022transfer}, jobBERT-de used a 20-epoch MaChAmp configuration and ESCOXLM-R used its released five-epoch configuration~\citep{zhang2023escoxlmr}. Both used five recorded seeds. Compatibility fields were added for newer MaChAmp without changing the learning rate or epoch limit; the exact changes are archived. The runs therefore use the public releases with compatibility changes, rather than recreating every original environment. They contain 2,557 test sentences, whereas the ESCOXLM-R paper reports 2,943. The ICT extraction task differs from the EDU/EXP/LNG extraction and taxonomy-classification pipeline described by \citet{gnehm2022skills}.

\subsection{Data and Implementation Sources}
We obtained the common-protocol data from the public Hugging Face releases of
\href{https://huggingface.co/datasets/jjzha/skillspan}{SkillSpan},
\href{https://huggingface.co/datasets/jjzha/kompetencer}{Kompetencer},
\href{https://huggingface.co/datasets/jjzha/gnehm}{Gnehm ICT},
\href{https://huggingface.co/datasets/jjzha/green}{Green},
\href{https://huggingface.co/datasets/jjzha/sayfullina}{Sayfullina}, and
\href{https://huggingface.co/datasets/jjzha/fijo}{FIJO}.
The corresponding papers are cited in the evaluation protocol; preprocessing and splits in these releases may differ from the original studies. The common protocol uses released BIO extraction streams: Sayfullina is a derived extraction task rather than its original candidate-phrase disambiguation task, and Kompetencer extraction is separate from the fine-grained classification of supplied spans reported below. Code is available from the
\href{https://github.com/kris927b/SkillSpan}{SkillSpan},
\href{https://github.com/jjzha/kompetencer}{Kompetencer},
\href{https://github.com/mainlp/escoxlmr}{ESCOXLM-R}, and
\href{https://github.com/machamp-nlp/machamp}{MaChAmp} repositories.
The \href{https://github.com/jjzha/nnose}{NNOSE repository} supplies the retrieval-augmented implementation used in the subsequent A100 experiments.
Our linear-head protocol and compatibility changes are described above. Recorded revisions are listed in the experiment records below.

\subsection{Uncertainty and Experiment Records}
Table~\ref{tab:external-common-ci} gives the paired sentence-bootstrap intervals. Sayfullina's interval includes zero, although its lower endpoint is close to zero. Because the intervals are pointwise, exclusions of zero do not constitute significance tests adjusted for multiple comparisons. The intervals capture sentence sampling uncertainty conditional on the trained seeds; seed SD describes variation across runs. Neither accounts for dependence between sentences from the same advertisement.

The \href{https://github.com/AlfredJamesLi/chinese-skillspan-benchmark/blob/36f009fab01b8433355cb7ef5af8fad02c6cc519/reproduction/process_archive_20260924/README.md}{versioned experiment records} link the common-protocol, Chinese adaptation, public-subset extraction, Kompetencer classification, and NNOSE retrieval experiments. They retain all seeds, model revisions, selected epochs, per-seed scores, and the available configuration and prediction hashes. Some dataset and executed-code identifiers remain unresolved.

External-task checks reconciled per-seed summaries with the ledgers; they did not rerun training or independently rescore predictions. For the Chinese encoders, local \texttt{cnss-lskt-1.2.0} rescoring reproduced all six runs' exact, relaxed, boundary-only, and per-type scores. Development-selected epochs and source offsets were also checked. Server-reported checkpoint re-inference matched archived tags on all 150 sentences per run; this was not a training rerun. The \href{https://github.com/AlfredJamesLi/chinese-skillspan-benchmark/tree/0c73207c0f9b00fa87f385b78bcb7342bd00bfcb/results_snapshots/table11_evidence_20260924}{pinned audit archive} and \href{https://github.com/AlfredJamesLi/chinese-skillspan-benchmark/blob/main/reproduction/evidence_review_20260925/README.md}{verification record} preserve code, configurations, predictions, and the detailed checks.

\begin{table}[H]
\centering\small
\caption{Paired differences under the common protocol (ESCOXLM-R minus XLM-R).}
\label{tab:external-common-ci}
\begin{tabularx}{\linewidth}{@{}L P{.19\linewidth}rr@{}}
\toprule
Dataset & Stream & Mean difference & 95\% interval \\
\midrule
SkillSpan & Skill & $+0.010$ & $[-0.003, 0.023]$ \\
SkillSpan & Knowledge & $-0.027$ & $[-0.041, -0.013]$ \\
Kompetencer & Skill & $-0.018$ & $[-0.063, 0.022]$ \\
Kompetencer & Knowledge & $-0.025$ & $[-0.102, 0.056]$ \\
Gnehm ICT & ICT & $-0.023$ & $[-0.034, -0.013]$ \\
Green & Native types & $-0.001$ & $[-0.021, 0.017]$ \\
Sayfullina & Skill & $+0.006$ & $[-1.47\times10^{-5}, 0.011]$ \\
FIJO & Native types & $-0.004$ & $[-0.043, 0.038]$ \\
\bottomrule
\end{tabularx}
\par\smallskip\RaggedRight\footnotesize Pointwise sentence-bootstrap intervals (2,000 resamples; seed 20260921), conditional on seeds 42/43/44. No advertisement clustering or multiple-comparison adjustment. Values are rounded after computing differences and intervals. Three decimal places are used except for Sayfullina's near-zero lower endpoint, shown with three significant digits in scientific notation to preserve its negative sign. Its interval includes zero.
\end{table}

\begin{table}[H]
\centering\small
\caption{SkillSpan test results: published single-task references and MaChAmp reruns on the public HOUSE and TECH subsets (0--1; five seeds).}
\label{tab:external-native}
\begin{tabularx}{\linewidth}{@{}L P{.15\linewidth}rrr@{}}
\toprule
Model & Task & Published TEST & Rerun HOUSE & Rerun TECH \\
\midrule
BERT & Skills & $0.543\pm0.007$ & $0.496\pm0.013$ & $0.514\pm0.012$ \\
BERT & Knowledge & $0.624\pm0.004$ & $0.550\pm0.004$ & $0.700\pm0.008$ \\
SpanBERT & Skills & --- & $0.516\pm0.014$ & $0.526\pm0.014$ \\
SpanBERT & Knowledge & --- & $0.536\pm0.017$ & $0.691\pm0.015$ \\
JobBERT & Skills & $0.561\pm0.005$ & $0.534\pm0.014$ & $0.538\pm0.013$ \\
JobBERT & Knowledge & $0.639\pm0.003$ & $0.577\pm0.004$ & $0.708\pm0.005$ \\
JobSpanBERT & Skills & $0.566\pm0.008$ & $0.543\pm0.013$ & $0.519\pm0.010$ \\
JobSpanBERT & Knowledge & $0.611\pm0.010$ & $0.546\pm0.018$ & $0.687\pm0.014$ \\
\bottomrule
\end{tabularx}
\par\smallskip\RaggedRight\footnotesize Published values are the TEST/STL results in Table 5 of \citet{zhang2022skillspan}, converted from percent and rounded to three decimals; their HOUSE/TECH rows report development results. Reruns use MaChAmp per-site span-F1 on public HOUSE+TECH, excluding BIG. The evaluation samples differ, so score differences cannot be interpreted as matched replication gaps. The source table gives no SpanBERT TEST score. Multi-task training was not rerun.
\end{table}


\begin{table}[H]
\centering\small
\caption{German public ICT extraction: native-task recipes (0--1; five seeds).}
\label{tab:external-gnehm}
\begin{tabularx}{\linewidth}{@{}L P{.25\linewidth}rr@{}}
\toprule
Task/model & Recipe & Published & Rerun \\
\midrule
Public ICT: jobBERT-de & 20-epoch MaChAmp & --- & $0.885\pm0.010$ \\
Public ICT: ESCOXLM-R & 5-epoch released recipe & $0.884\pm0.005$ & $0.891\pm0.005$ \\
\bottomrule
\end{tabularx}
\par\smallskip\RaggedRight\footnotesize Reruns use 2,557 test sentences. The published ESCOXLM-R score comes from Table 2 of \citet{zhang2023escoxlmr}, whose Table 1 reports 2,943 test sentences. These different releases prevent a matched comparison with the published score. No corresponding published jobBERT-de score is supplied. Unlike the common-protocol comparison, these two native-task configurations use different training budgets measured in epochs.
\end{table}


\FloatBarrier
\subsection{Kompetencer Classification and NNOSE Reruns}
Kompetencer fine-grained classification~\citep{jensen2022kompetencer} and NNOSE~\citep{zhang2024nnose} were run on an NVIDIA A100 (80 GB). Software versions and compatibility records are provided in the \href{https://github.com/AlfredJamesLi/chinese-skillspan-benchmark/blob/1d1c24e20ee049fc768034b5baf58fa882e08e2c/reproduction/process_archive_20260924/secondary_review/source/tex/external_benchmarks/paper_A_native.tex}{native-task configurations}. These experiments retain their task-specific metrics and are separate from the common linear-head protocol.

For Kompetencer, we compared DaJobBERT with RemBERT~\citep{chung2021rembert}. MaChAmp used a 20-epoch budget and the five seeds recorded for the SkillSpan native runs. Training used Danish (DA), English (EN), or both (EN+DA), with 138 Danish and 9,472 English training records. All conditions selected checkpoints on the 1,577-record English development set. Table~\ref{tab:external-kompetencer-classification} reports unweighted macro-F1 for classification of supplied spans. The recorded DaJobBERT initialization is \href{https://github.com/AlfredJamesLi/chinese-skillspan-benchmark/blob/b84f218b7142c97ab03a60499a9b30e7ec332cab/results_snapshots/a_native_20260924/A_NATIVE_RECEIPT.md}{\texttt{jjzha/dajobbert-base-uncased}}; the linked receipt retains its full revision. Together with the unweighted metric, this makes the experiment a task-specific rerun rather than a matched reproduction of the original cased setting.

NNOSE used JobBERT and JobBERTa with seed 113412 and development-based early stopping. Each datastore contained only training examples from the evaluated dataset. Neighbour count $k$, interpolation weight $\lambda$, and temperature $T$ were selected on development data and frozen before testing. Table~\ref{tab:external-nnose} compares seqeval span-F1 from the linear head and retrieval-augmented predictions of the same checkpoint. These linear-head results belong to the NNOSE pipeline, not the common-protocol XLM-R comparison. The all-dataset datastore was not run, and the single-seed results do not establish a general retrieval benefit.

\begin{table}[htbp]
\centering\small
\caption{Kompetencer fine-grained classification: macro-F1 across five seeds (0--1).}
\label{tab:external-kompetencer-classification}
\begin{tabularx}{\linewidth}{@{}L P{.18\linewidth}rr@{}}
\toprule
Encoder & Training & Danish test & English test \\
\midrule
DaJobBERT & DA & $0.182\pm0.034$ & $0.064\pm0.012$ \\
DaJobBERT & EN & $0.232\pm0.028$ & $0.539\pm0.010$ \\
DaJobBERT & EN+DA & $0.340\pm0.017$ & $0.529\pm0.008$ \\
RemBERT & DA & $0.050\pm0.072$ & $0.034\pm0.041$ \\
RemBERT & EN & $0.261\pm0.019$ & $0.552\pm0.017$ \\
RemBERT & EN+DA & $0.316\pm0.019$ & $0.554\pm0.014$ \\
\bottomrule
\end{tabularx}
\par\smallskip\RaggedRight\footnotesize Mean $\pm$ sample SD; Danish/English test sizes are 784/1,578 classified spans. All conditions use English development data. EN training with Danish testing is the EN-to-DA transfer condition. DaJobBERT resolved to an uncased checkpoint despite the requested cased identifier. Scores are unweighted classification macro-F1, not span-F1 or the source plot's weighted metric. No matched published-score comparison is claimed.
\end{table}

\begin{table}[htbp]
\centering\small
\caption{NNOSE with an in-dataset training datastore: test span-F1 (0--1).}
\label{tab:external-nnose}
\begin{tabularx}{\linewidth}{@{}L P{.22\linewidth}rrr@{}}
\toprule
Dataset & Encoder & Test $n$ & Linear & kNN \\
\midrule
SkillSpan & JobBERT & 3,569 & 0.613 & 0.612 \\
SkillSpan & JobBERTa & 3,569 & 0.631 & 0.625 \\
Green & JobBERT & 335 & 0.491 & 0.500 \\
Green & JobBERTa & 335 & 0.466 & 0.485 \\
Sayfullina & JobBERT & 1,851 & 0.897 & 0.896 \\
Sayfullina & JobBERTa & 1,851 & 0.915 & 0.912 \\
\bottomrule
\end{tabularx}
\par\smallskip\RaggedRight\footnotesize Single seed 113412; no SD is estimated. Selected epochs and development-selected retrieval settings are preserved in the \href{https://github.com/AlfredJamesLi/chinese-skillspan-benchmark/blob/1d1c24e20ee049fc768034b5baf58fa882e08e2c/reproduction/process_archive_20260924/secondary_review/tables/nnose_full.tex}{complete run table}. Both columns use the same trained checkpoint within each row. The datastore excludes development and test examples. SkillSpan uses public HOUSE+TECH. The all-dataset datastore was not evaluated, and alignment with an original-paper result cell has not been established.
\end{table}

\FloatBarrier

Combined English/Danish training improved Danish macro-F1 for both encoders; RemBERT's Danish-only score varied across seeds ($0.050\pm0.072$). English-based selection and the unweighted metric prevent a matched published-score comparison. NNOSE improved Green but slightly reduced SkillSpan and Sayfullina F1 in these single-seed runs, offering no consistent retrieval benefit.

The \href{https://github.com/AlfredJamesLi/chinese-skillspan-benchmark/blob/1d1c24e20ee049fc768034b5baf58fa882e08e2c/reproduction/process_archive_20260924/secondary_review/source/tex/external_benchmarks/paper_A_native.tex}{coverage record} identifies uncompleted configurations; no missing result is assigned a zero score.

\FloatBarrier
